\begin{abstract}
This paper presents the design and implementation of a hybrid peer-to-peer (P2P) chat system, developed from first principles to satisfy the rigorous constraints of an undergraduate computer networking assignment. The core challenge was to create a real-time messaging application that integrates both Client-Server and P2P communication paradigms without the use of external frameworks or modern real-time protocols such as WebSockets. To overcome the browser's inherent inability to establish raw TCP connections, a novel hybrid architecture was devised. This architecture features a centralized backend server for user authentication, session management, and peer discovery, operating on a Client-Server model via HTTP. For real-time messaging, a decentralized network of P2P daemons, one for each user, establishes direct TCP socket connections for low-latency data exchange. The system employs an "HTTP Bridge" pattern, where the backend server acts as an intermediary, translating HTTP requests from the browser into commands for the corresponding P2P daemon. Real-time event and message delivery to the client is simulated using periodic HTTP polling (short-poll); the server responds immediately and the browser polls at regular intervals. The entire system is implemented in Python, relying exclusively on standard libraries such as socket and threading to handle low-level networking and concurrency. The resulting application is a robust, multi-process system that successfully demonstrates the practical integration of complex networking concepts and fulfills all academic requirements.
\end{abstract}